\subsection{Quartic potential identities}
\label{app:quartic-potential-details}

For the normalization of \cref{ex:Riedler-correct-power}, take
$G=RQ^{3/4}F$ and use the differential table
\eqref{eq:quartic-ambient-closure}.
The quotient chain rule gives
\begin{align}
 |\nabla G|^2
 &=RQ^{1/2}\bigl(16R^2Q^2+33R^2F^2+32QF^2\bigr),
 \label{eq:Riedler-correct-gradient}\\
 \frac{\cL(G)}{F}
 &=RQ^{1/4}\Bigl[
 16(12k-47)R^4Q^2+16(2N-41)R^2Q^3\notag\\
 &\qquad
 +F^2\bigl((396k+66)R^4+(384k+66N-681)R^2Q\notag\\
 &\hspace{35mm} +(64N+160)Q^2\bigr)\Bigr].
 \label{eq:Riedler-correct-defect}
\end{align}
Every coefficient is positive for $N\geq32$ and $k\geq4$.
Equivalently, apply \cref{prop:combined-power-multiplier} with
\eqref{eq:Riedler-correct-potential}.


On the regular zero set, these formulas yield
\eqref{eq:Riedler-positive-boundary}.

\paragraph*{The $(3,7,7)$ specialization}
The differential table and certificate used for
\cref{thm:weighted-Clifford-quartic} are
\cref{eq:balanced-quartic-table,eq:balanced-quartic-certificate}
at $N=32$ and $\nu=7$:
\begin{equation}\label{eq:weighted-Clifford-table}
 \begin{aligned}
 \Delta(R,W,F)&=(64,28R,0),&\quad \Gamma(F,W)&=4RF,\\
 \Gamma(F,F)&=\Gamma(W,W)=4RW.
 \end{aligned}
\end{equation}
\begin{equation}\label{eq:weighted-Clifford-certificate}
 \cL(WF)=16WF\left[(3R^2-2W)W^2+(21R^2-6W)F^2\right].
\end{equation}
The certificate file is
\nolinkurl{LSC-WEIGHTED-CLIFFORD-QUARTIC.wl}.
